% 第29批 40题及解析候选；未达到高级综合题准入门槛，未独立复核。

\section*{第29批：40题候选题解}

\subsection*{M193\quad 圆盘径向滑珠}

匀速转盘径向无摩擦槽，角速率$\Omega$，滑珠初始位置$r_0$，初始径向速度零。

\par\textbf{分问：} (1) 径向方程；(2) 位移时间式；(3) 槽壁横向力；(4) 旋转驱动功率。

\par\textbf{解析：} (1) $\ddot r=\Omega^2r$；(2) $r=r_0\cosh\Omega t$；(3) $F_\theta=2m\Omega\dot r$；(4) $P=2m\Omega^2r\dot r$。

\medskip

\subsection*{M194\quad 球面脱离}

小球从光滑球面顶端静止出发，在微扰下滑，球半径$R$。

\par\textbf{分问：} (1) 速度与角度；(2) 法向约束；(3) 脱离角；(4) 脱离速度。

\par\textbf{解析：} (1) $v^2=2gR(1-\cos\theta)$；(2) $N=mg\cos\theta-mv^2/R$；(3) $\cos\theta=2/3$；(4) $v=\sqrt{2gR/3}$。

\medskip

\subsection*{M195\quad 均匀行星隧道}

半径$R$质量$M$的均匀球体沿直径开光滑隧道。

\par\textbf{分问：} (1) 内部力；(2) 角频率；(3) 穿越时间；(4) 中心速度。

\par\textbf{解析：} (1) $F=-GMmx/R^3$；(2) $\omega=\sqrt{GM/R^3}$；(3) $t=\pi/\omega$；(4) $v_c=\omega R$。

\medskip

\subsection*{M196\quad 有惯量阿特伍德机}

两重物$m_1>m_2$经不滑动轻绳跨过转动惯量$I$、半径$R$的轮。

\par\textbf{分问：} (1) 加速度；(2) 两侧绳张力；(3) 角加速度；(4) 从静止下降$h$后的速度。

\par\textbf{解析：} (1) $a=(m_1-m_2)g/(m_1+m_2+I/R^2)$；(2) $T_1=m_1(g-a),\ T_2=m_2(g+a)$；(3) $\alpha=a/R$；(4) $v=\sqrt{2ah}$。

\medskip

\subsection*{M197\quad 反平方中心力}

粒子质量$m$，径向吸引力$-k/r^2$，角动量大小$L$。

\par\textbf{分问：} (1) 角动量守恒；(2) 有效势；(3) 圆轨道；(4) 稳定性。

\par\textbf{解析：} (1) $\dot L=0$；(2) $U_{\rm eff}=L^2/(2mr^2)-k/r$；(3) $r_c=L^2/(mk)$；(4) $U_{\rm eff}''(r_c)=k/r_c^3>0$。

\medskip

\subsection*{M198\quad 两块滑动摩擦}

上块$m$初速$v_0$，下板$M$静止，地面光滑，两块间动摩擦$\mu$，直到共速。

\par\textbf{分问：} (1) 各自加速度；(2) 相对加速度；(3) 追及时间；(4) 共同速度。

\par\textbf{解析：} (1) $a_m=-\mu g,\ a_M=\mu mg/M$；(2) $a_{rel}=-\mu g(1+m/M)$；(3) $t=v_0/[\mu g(1+m/M)]$；(4) $v=mv_0/(m+M)$。

\medskip

\subsection*{M199\quad 径向引力抛射}

从距天体中心$r_0$处以径向速率$v_0$向外抛射，天体参数$GM$。

\par\textbf{分问：} (1) 比能量；(2) 速度半径关系；(3) 逃逸速度；(4) 未逃逸最大半径。

\par\textbf{解析：} (1) $\varepsilon=v_0^2/2-GM/r_0$；(2) $v_r^2=v_0^2+2GM(1/r-1/r_0)$；(3) $v_e=\sqrt{2GM/r_0}$；(4) $r_{max}=[1/r_0-v_0^2/(2GM)]^{-1}$（末式仅限未逃逸）。

\medskip

\subsection*{E193\quad 均匀带电实心球}

非导体球半径$R$，电荷$Q$均匀分布。

\par\textbf{分问：} (1) 体密度；(2) 内部电场；(3) 外部电场；(4) 中心表面电势差。

\par\textbf{解析：} (1) $\rho=3Q/(4\pi R^3)$；(2) $E_{in}=Qr/(4\pi\epsilon_0R^3)$；(3) $E_{out}=Q/(4\pi\epsilon_0r^2)$；(4) $\varphi(0)-\varphi(R)=Q/(8\pi\epsilon_0R)$。

\medskip

\subsection*{E194\quad 均匀介质平板电容器}

无限平行金属板间距$d$，自由面电荷$\pm\sigma$，板间介电常数$\epsilon$。

\par\textbf{分问：} (1) 位移矢量；(2) 电场；(3) 板间电压；(4) 单位面积储能。

\par\textbf{解析：} (1) $D=\sigma$；(2) $E=\sigma/\epsilon$；(3) $V=\sigma d/\epsilon$；(4) $U/A=\sigma^2d/(2\epsilon)$。

\medskip

\subsection*{E195\quad 粒子回旋运动}

质量$m$、电荷$q\ne0$的粒子垂直于匀强磁场$B$以速率$v$运动。

\par\textbf{分问：} (1) 半径；(2) 角频率；(3) 周期；(4) 动能四倍后的半径。

\par\textbf{解析：} (1) $r=mv/(|q|B)$；(2) $\omega=|q|B/m$；(3) $T=2\pi m/(|q|B)$；(4) $r_{new}=2r$（非相对论）。

\medskip

\subsection*{E196\quad RL电路起动}

直流电源$V$接串联$R,L$，初始电流零。

\par\textbf{分问：} (1) 微分方程；(2) 电流函数；(3) 电感端电压；(4) 最终磁能。

\par\textbf{解析：} (1) $L\dot i+Ri=V$；(2) $i=V(1-e^{-Rt/L})/R$；(3) $V_L=Ve^{-Rt/L}$；(4) $U_L=LV^2/(2R^2)$。

\medskip

\subsection*{E197\quad 直导线框磁通}

电流$I(t)$无限长导线与矩形线框共面，线框长$\ell$，近远边距离$a,b$。

\par\textbf{分问：} (1) 磁场；(2) 磁通；(3) 感应电动势；(4) 电流增大时楞次方向。

\par\textbf{解析：} (1) $B=\mu_0I/(2\pi r)$；(2) $\Phi=\mu_0I\ell\ln(b/a)/(2\pi)$；(3) $\mathcal E=-\mu_0\ell\ln(b/a)\dot I/(2\pi)$；(4) 感应磁场抵抗穿框磁通增加；具体顺逆时针依观察法线。

\medskip

\subsection*{E198\quad 孤立金属球电容}

孤立金属球半径$R$带电$Q$，取无穷远势零。

\par\textbf{分问：} (1) 球外电势；(2) 表面电势；(3) 电容；(4) 静电能。

\par\textbf{解析：} (1) $\varphi(r)=Q/(4\pi\epsilon_0r)$；(2) $V=Q/(4\pi\epsilon_0R)$；(3) $C=4\pi\epsilon_0R$；(4) $U=Q^2/(8\pi\epsilon_0R)$。

\medskip

\subsection*{E199\quad 平行电流作用}

相距$d$的两根无限长平行导线电流$I_1,I_2$同向。

\par\textbf{分问：} (1) 第一根磁场；(2) 第二根单位长度受力；(3) 方向；(4) 距离加倍结果。

\par\textbf{解析：} (1) $B_1=\mu_0I_1/(2\pi d)$；(2) $F/\ell=\mu_0I_1I_2/(2\pi d)$；(3) 同向相吸；(4) 单位长度力减半。

\medskip

\subsection*{T186\quad 等压加热}

$n$摩尔理想气体，摩尔热容$C_p$常量，温差$\Delta T$。

\par\textbf{分问：} (1) 体积变化；(2) 气体做功；(3) 内能变化；(4) 吸热。

\par\textbf{解析：} (1) $\Delta V=nR\Delta T/p$；(2) $W=nR\Delta T$；(3) $\Delta U=n(C_p-R)\Delta T$；(4) $Q=nC_p\Delta T$。

\medskip

\subsection*{T187\quad 有限热源可逆取功}

两相同热容$C$物体温度$T_h>T_c$，绝热系统内用可逆热机使其达共同温度。

\par\textbf{分问：} (1) 总熵关系；(2) 最终温度；(3) 最大功；(4) 直接接触末温。

\par\textbf{解析：} (1) $\ln(T_f/T_h)+\ln(T_f/T_c)=0$；(2) $T_f=\sqrt{T_hT_c}$；(3) $W=C(T_h+T_c-2\sqrt{T_hT_c})$；(4) $T_{direct}=(T_h+T_c)/2$。

\medskip

\subsection*{T188\quad 多方气体过程}

理想气体遵从$pV^s=K$，$s\ne1$，从$(p_1,V_1)$到$V_2$。

\par\textbf{分问：} (1) 末压；(2) 边界功；(3) 终温比；(4) 热量。

\par\textbf{解析：} (1) $p_2=p_1(V_1/V_2)^s$；(2) $W=(p_2V_2-p_1V_1)/(1-s)$；(3) $T_2/T_1=(V_2/V_1)^{1-s}$；(4) $Q=nC_V(T_2-T_1)+W$。

\medskip

\subsection*{T189\quad 辐射散热}

热容$C$物体，面积$A$发射率$\varepsilon$，环境温度$T_a$。

\par\textbf{分问：} (1) 净辐射功率；(2) 冷却方程；(3) 小温差近似；(4) 弛豫时间。

\par\textbf{解析：} (1) $P=\varepsilon\sigma A(T^4-T_a^4)$；(2) $C\dot T=-P$；(3) $P\simeq4\varepsilon\sigma AT_a^3(T-T_a)$；(4) $\tau=C/(4\varepsilon\sigma AT_a^3)$。

\medskip

\subsection*{T190\quad 等容理想气体熵}

$n$摩尔定热容$C_V$气体等容升温$T_1$到$T_2$。

\par\textbf{分问：} (1) 边界功；(2) 热量；(3) 熵变；(4) 自由能微分恒等式。

\par\textbf{解析：} (1) $W=0$；(2) $Q=nC_V(T_2-T_1)$；(3) $\Delta S=nC_V\ln(T_2/T_1)$；(4) $dF=-S\,dT-p\,dV$，等容时$dF=-S\,dT$。

\medskip

\subsection*{T191\quad 跨热阻传热熵产}

两热库$T_h>T_c$经恒热阻$R_{th}$传热。

\par\textbf{分问：} (1) 热流；(2) 高温库熵率；(3) 低温库熵率；(4) 总熵产率。

\par\textbf{解析：} (1) $\dot Q=(T_h-T_c)/R_{th}$；(2) $\dot S_h=-\dot Q/T_h$；(3) $\dot S_c=\dot Q/T_c$；(4) $\dot S_{gen}=\dot Q(1/T_c-1/T_h)>0$。

\medskip

\subsection*{T192\quad 绝热自由膨胀}

理想气体绝热向真空自由膨胀$V_i\to V_f>V_i$。

\par\textbf{分问：} (1) 对外做功；(2) 热量；(3) 温变；(4) 熵变。

\par\textbf{解析：} (1) $W=0$；(2) $Q=0$；(3) $\Delta T=0$（理想气体内能仅依温度）；(4) $\Delta S=nR\ln(V_f/V_i)>0$。

\medskip

\subsection*{W186\quad 单原子晶格色散}

无限质量$m$、间距$a$、相邻耦合弹簧$K$的链。

\par\textbf{分问：} (1) 运动方程；(2) 色散关系；(3) 长波声速；(4) 区边界频率。

\par\textbf{解析：} (1) $m\ddot u_j=K(u_{j+1}-2u_j+u_{j-1})$；(2) $\omega^2=4K\sin^2(ka/2)/m$；(3) $c=a\sqrt{K/m}$；(4) $\omega(\pi/a)=2\sqrt{K/m}$。

\medskip

\subsection*{W187\quad 简谐振子守恒量}

位移$x=A\cos(\omega t+\phi)$，弹性系数$k=m\omega^2$。

\par\textbf{分问：} (1) 速度；(2) 动能；(3) 势能；(4) 总能量。

\par\textbf{解析：} (1) $v=-A\omega\sin(\omega t+\phi)$；(2) $K=mA^2\omega^2\sin^2(\omega t+\phi)/2$；(3) $U=kA^2\cos^2(\omega t+\phi)/2$；(4) $E=kA^2/2$。

\medskip

\subsection*{W188\quad 两弦接头反射}

张力同为$F$，两侧波速$v_1,v_2$，从1侧正入射位移波。

\par\textbf{分问：} (1) 位移连续；(2) 斜率连续；(3) 位移反射率；(4) 位移透射率。

\par\textbf{解析：} (1) $A_i+A_r=A_t$；(2) $k_1(A_i-A_r)=k_2A_t$；(3) $r=(v_2-v_1)/(v_2+v_1)$；(4) $t=2v_2/(v_1+v_2)$。

\medskip

\subsection*{W189\quad 波的端点反射}

半无限弦位于$x<0$，横波在$x=0$端反射。

\par\textbf{分问：} (1) 固定端边界；(2) 固定端反射相位；(3) 自由端边界；(4) 自由端反射相位。

\par\textbf{解析：} (1) $y(0,t)=0$；(2) 反射位移振幅负号，差相$\pi$；(3) $\partial_xy(0,t)=0$；(4) 反射位移振幅正号，差相零。

\medskip

\subsection*{W190\quad 高斯包色散}

$\omega=\omega_0+v_g(k-k_0)+\beta_2(k-k_0)^2/2$，初始无啁啾高斯包标准差$\sigma_0$。

\par\textbf{分问：} (1) 中心运动；(2) 初始波数宽；(3) 时变空间宽；(4) 不展宽条件。

\par\textbf{解析：} (1) $x_c(t)=x_c(0)+v_gt$；(2) $\sigma_k=1/(2\sigma_0)$；(3) $\sigma_x=\sigma_0\sqrt{1+[\beta_2t/(2\sigma_0^2)]^2}$；(4) $\beta_2=0$（二阶近似）。

\medskip

\subsection*{W191\quad 开口管驻波声能}

长度$L$、截面积$S$的双开口管，压力驻波$p'=p_a\sin(n\pi x/L)\cos\omega t$。

\par\textbf{分问：} (1) 波数；(2) 频率；(3) 压力节点数；(4) 平均压缩势能。

\par\textbf{解析：} (1) $k=n\pi/L$；(2) $f=nc/(2L)$；(3) $n+1$（含两端）；(4) $\overline U_p=SLp_a^2/(8\rho c^2)$。

\medskip

\subsection*{W192\quad 混合边界弦驻波}

弦长$L$，一端固支、另一端自由，波速$v$。

\par\textbf{分问：} (1) 本征函数；(2) 边界波数；(3) 频率谱；(4) 次低频与基频之比。

\par\textbf{解析：} (1) $X(x)=A\sin kx$；(2) $k_n=(n+1/2)\pi/L$；(3) $f_n=(2n+1)v/(4L)$，$n\ge0$；(4) $f_1/f_0=3$。

\medskip

\subsection*{O174\quad 垂直入射菲涅耳公式}

光在非磁性无损两介质$n_1,n_2$界面正入射。

\par\textbf{分问：} (1) 电场反射系数；(2) 电场透射系数；(3) 反射率；(4) 透射率及守恒。

\par\textbf{解析：} (1) $r=(n_1-n_2)/(n_1+n_2)$；(2) $t=2n_1/(n_1+n_2)$；(3) $R=r^2$；(4) $T=4n_1n_2/(n_1+n_2)^2$，$R+T=1$。

\medskip

\subsection*{O175\quad 正入射薄膜干涉}

空气-薄膜-低折射率基底，薄膜$n>1$、厚度$d$，低反射率双光束近似。

\par\textbf{分问：} (1) 往返光程；(2) 反射相移；(3) 相长条件；(4) 最薄非零相长厚度。

\par\textbf{解析：} (1) $2nd$；(2) 空气-膜反射发生$\pi$相移，膜-低折射率基底不发生；(3) $2nd=(m+1/2)\lambda_0$；(4) $d_{min}=\lambda_0/(4n)$。

\medskip

\subsection*{O176\quad 两路干涉仪}

理想等振幅马赫-曾德尔干涉仪，两臂光程差$\Delta\ell$，波长$\lambda$。

\par\textbf{分问：} (1) 相位差；(2) 第一输出功率；(3) 第二输出功率；(4) 总功率。

\par\textbf{解析：} (1) $\delta=2\pi\Delta\ell/\lambda$；(2) $I_1=I_0\cos^2(\delta/2)$；(3) $I_2=I_0\sin^2(\delta/2)$；(4) $I_1+I_2=I_0$。

\medskip

\subsection*{O177\quad 对称棱镜色散}

空气中顶角$A$、折射率$n$的棱镜测得最小偏向$D_m$。

\par\textbf{分问：} (1) 折射角约束；(2) 对称光线条件；(3) 折射率；(4) 小角偏向。

\par\textbf{解析：} (1) $r_1+r_2=A$；(2) $i_1=i_2=(A+D_m)/2$；(3) $n=\sin[(A+D_m)/2]/\sin(A/2)$；(4) $D_m\simeq(n-1)A$（弧度小角）。

\medskip

\subsection*{O178\quad 圆形孔径衍射}

直径$D$圆孔照明，单色波长$\lambda$，后有焦距$f$透镜。

\par\textbf{分问：} (1) 第一暗环角度；(2) 焦面半径；(3) 瑞利分辨角；(4) 口径倍增后变化。

\par\textbf{解析：} (1) $\theta=1.22\lambda/D$；(2) $r=1.22\lambda f/D$；(3) $\theta_R=1.22\lambda/D$；(4) 角分辨阈值减半。

\medskip

\subsection*{O179\quad 偏振与分析器}

线偏振入射强度$I_0$，分析器轴夹角$\theta$。

\par\textbf{分问：} (1) 透射振幅；(2) 透射强度；(3) 垂直偏振结果；(4) 平行偏振结果。

\par\textbf{解析：} (1) $E_t=E_0\cos\theta$；(2) $I=I_0\cos^2\theta$；(3) $I(\pi/2)=0$；(4) $I(0)=I_0$。

\medskip

\subsection*{Q174\quad 相对论能动量}

质量$m$粒子速率$v<c$。

\par\textbf{分问：} (1) 总能量；(2) 动量；(3) 壳层关系；(4) 小速动能。

\par\textbf{解析：} (1) $E=\gamma mc^2$；(2) $p=\gamma mv$；(3) $E^2-p^2c^2=m^2c^4$；(4) $K\simeq mv^2/2+3mv^4/(8c^2)$。

\medskip

\subsection*{Q175\quad 电压加速电子波长}

电子静止经电压$V$加速，先假设非相对论。

\par\textbf{分问：} (1) 动能；(2) 动量；(3) 德布罗意波长；(4) 适用条件。

\par\textbf{解析：} (1) $K=eV$；(2) $p=\sqrt{2m_eeV}$；(3) $\lambda=h/\sqrt{2m_eeV}$；(4) $eV\ll m_ec^2$。

\medskip

\subsection*{Q176\quad 矩形量子势垒}

势垒$V_0>E>0$、宽$a$，质量$m$粒子一维散射。

\par\textbf{分问：} (1) 势垒内方程；(2) 衰减常数；(3) 通解；(4) 厚势垒透射主项。

\par\textbf{解析：} (1) $-\hbar^2\psi''/(2m)+(V_0-E)\psi=0$；(2) $\kappa=\sqrt{2m(V_0-E)}/\hbar$；(3) $\psi=Ae^{\kappa x}+Be^{-\kappa x}$；(4) $T\propto e^{-2\kappa a}$，$\kappa a\gg1$。

\medskip

\subsection*{Q177\quad 氢样离子光谱}

核电荷$Ze$、电子核约化质量$\mu$，从$n_i$跃迁至$n_f<n_i$。

\par\textbf{分问：} (1) 玻尔能级；(2) 发射能量；(3) 波长；(4) 核电荷标度。

\par\textbf{解析：} (1) $E_n=-C_Z/n^2$，$C_Z=\mu Z^2e^4/[2(4\pi\epsilon_0)^2\hbar^2]$；(2) $h\nu=C_Z(1/n_f^2-1/n_i^2)$；(3) $\lambda=hc/[C_Z(1/n_f^2-1/n_i^2)]$；(4) 固定约化质量时能量正比$Z^2$。

\medskip

\subsection*{Q178\quad 一维势阱光谱}

无限深势阱宽度$L$，质量$m$粒子。

\par\textbf{分问：} (1) 归一波函数；(2) 基态至第一激发态能差；(3) 光频；(4) 阱宽加倍影响。

\par\textbf{解析：} (1) $\psi_n=\sqrt{2/L}\sin(n\pi x/L)$；(2) $\Delta E=3\pi^2\hbar^2/(2mL^2)$；(3) $\nu=\Delta E/h$；(4) 频率变为$1/4$。

\medskip

\subsection*{Q179\quad 自旋二能级系统}

自旋$1/2$磁矩$\boldsymbol\mu=\gamma_s\mathbf S$，磁场$B\hat z$，温度$T$。

\par\textbf{分问：} (1) 哈密顿量；(2) 两个能级；(3) 劈裂；(4) 玻尔兹曼布居比。

\par\textbf{解析：} (1) $H=-\gamma_sBS_z$；(2) $E_\pm=\mp\gamma_s\hbar B/2$；(3) $\Delta E=|\gamma_s|\hbar B$；(4) $p_+/p_-=\exp[-(E_+-E_-)/(k_BT)]$。

\medskip
